\section{Future Work}
Although the examples shown show how valuable a tool like
Cedar can be, it also highlights several avenues for future development.

\subsection{Language and Compiler Features}
There are a few features that were considered
out of scope for this initial implementation, however
they heavily limited the real world examples that could be tested
with Cedar. Implementing these features would greatly increase
the usefulness of the language and set it up as a tool useful
for real world systems programming and research.

\begin{enumerate}
    \item Complete well-formedness checking: The current implementation
    does not perform a full static analysis of the layout.
    Implementing this should be a priority, and the current design
    leaves space for this to be added in future work. The CR stage
    calculates offsets and sizes required for this check to to take
    place, but it is not yet implemented.
    \item Missing Structures: Bit fields, unions and pointers
    are commonly found structures in C programs, and their support
    would immediately increase the amount of layouts that could be expressed
    with the language.
    \item Unified Array accessors: A trivial but useful addition
    would be to ensure the codegen stage generates a unified
    accessor for a whole array, as well as per-element accessors. 
    This would reduce the amount of glue code required to
    integrate Cedar intro existing C codebases.
\end{enumerate}

\subsection{Towards a Verified System}
As with its predecessors, Cedar was conceived as an idea
for a verified system for data-layouts.
The design of the language reflects this, and undertaking
the task of ensuring each of the internal translations
stages are formally verified, would be the first step into
turning Cedar into a verified tool that can integrate easily
with verified compilers such as CompCert in order to provide
end-to-end guarantees about code behavior.

\textbf{Formalizing Internal Representations}:

Each translation stage (L, LR and CR) encodes
a refinement of information. From a 
symbolic structure, to resolved offsets and finally concrete
bit ranges.
These transformations are functional and self-contained, enabling
automated reasoning about their correctness with a proof assistant such
as Roq/Coq.
Formalizing their behavior would enable proofs that layout meaning
is preserved across stages: that the semantics of a record array
in L correspond exactly to its bit-range representation in CR.
Such results would elevate Cedar into a verifiable component
with a broader compilation framework.

\textbf{Integrating with Verified Compilation}:

An interesting direction would be to develop a custom C compiler,
most likely adapting CompCert as it already provides a verified
pipeline, and integrate the Cedar compiler into that pipeline.
This would allow for Cedar code to be written in-line along with 
C code, creating a seamless experience for programmers who
need a fully verified toolchain for systems programming.

\textbf{Tooling visualization}:

As Cedar's analysis grows more sophisticated, corresponding improvements
in usability will become essential. Interactive visualization of layouts,
showing byte ranges, field names and endianness, could transform the compiler
into a teaching and debugging aid. 
